\question{10}{Considere un grafo representado por las siguientes listas de adyacencia:

\begin{center}
\begin{tikzpicture}[baseline=(current bounding box.center),
  idxcell/.style={draw, minimum size=0.42cm, inner sep=0pt, outer sep=0pt},
  slot/.style={draw, minimum size=0.42cm, inner sep=1pt, outer sep=0pt, font=\scriptsize},
  -latex, thick,
]
  \node[idxcell]                    (idx0) at (0, 0) {};
  \node[idxcell, below=0pt of idx0] (idx1) {};
  \node[idxcell, below=0pt of idx1] (idx2) {};
  \node[idxcell, below=0pt of idx2] (idx3) {};
  \node[idxcell, below=0pt of idx3] (idx4) {};
  \node[idxcell, below=0pt of idx4] (idx5) {};
  \node[idxcell, below=0pt of idx5] (idx6) {};

  \node[left=6pt of idx0] {$1$};
  \node[left=6pt of idx1] {$2$};
  \node[left=6pt of idx2] {$3$};
  \node[left=6pt of idx3] {$4$};
  \node[left=6pt of idx4] {$5$};
  \node[left=6pt of idx5] {$6$};
  \node[left=6pt of idx6] {$7$};

  % idx0 (v1): 4, 6, 7
  \node[slot, right=0.35cm of idx0] (r0a1) {4};
  \node[slot, right=0pt of r0a1]    (r0a2) {};
  \node[slot, right=0.35cm of r0a2] (r0b1) {6};
  \node[slot, right=0pt of r0b1]    (r0b2) {};
  \node[slot, right=0.35cm of r0b2] (r0c1) {7};
  \node[slot, right=0pt of r0c1]    (r0c2) {};
  \draw (idx0.center) -- (r0a1);
  \draw (r0a2.center) -- (r0b1);
  \draw (r0b2.center) -- (r0c1);

  % idx1 (v2): 5, 6
  \node[slot, right=0.35cm of idx1] (r1a1) {5};
  \node[slot, right=0pt of r1a1]    (r1a2) {};
  \node[slot, right=0.35cm of r1a2] (r1b1) {6};
  \node[slot, right=0pt of r1b1]    (r1b2) {};
  \draw (idx1.center) -- (r1a1);
  \draw (r1a2.center) -- (r1b1);

  % idx2 (v3): 5, 7
  \node[slot, right=0.35cm of idx2] (r2a1) {5};
  \node[slot, right=0pt of r2a1]    (r2a2) {};
  \node[slot, right=0.35cm of r2a2] (r2b1) {7};
  \node[slot, right=0pt of r2b1]    (r2b2) {};
  \draw (idx2.center) -- (r2a1);
  \draw (r2a2.center) -- (r2b1);

  % idx3 (v4): 1
  \node[slot, right=0.35cm of idx3] (r3a1) {1};
  \node[slot, right=0pt of r3a1]    (r3a2) {};
  \draw (idx3.center) -- (r3a1);

  % idx4 (v5): 2, 3
  \node[slot, right=0.35cm of idx4] (r4a1) {2};
  \node[slot, right=0pt of r4a1]    (r4a2) {};
  \node[slot, right=0.35cm of r4a2] (r4b1) {3};
  \node[slot, right=0pt of r4b1]    (r4b2) {};
  \draw (idx4.center) -- (r4a1);
  \draw (r4a2.center) -- (r4b1);

  % idx5 (v6): 1, 2
  \node[slot, right=0.35cm of idx5] (r5a1) {1};
  \node[slot, right=0pt of r5a1]    (r5a2) {};
  \node[slot, right=0.35cm of r5a2] (r5b1) {2};
  \node[slot, right=0pt of r5b1]    (r5b2) {};
  \draw (idx5.center) -- (r5a1);
  \draw (r5a2.center) -- (r5b1);

  % idx6 (v7): 1, 3
  \node[slot, right=0.35cm of idx6] (r6a1) {1};
  \node[slot, right=0pt of r6a1]    (r6a2) {};
  \node[slot, right=0.35cm of r6a2] (r6b1) {3};
  \node[slot, right=0pt of r6b1]    (r6b2) {};
  \draw (idx6.center) -- (r6a1);
  \draw (r6a2.center) -- (r6b1);
\end{tikzpicture}
\end{center}

El primer elemento del arreglo de apuntadores representa el vértice \(v_1\). Cada vez que visite los vecinos de un vértice, hágalo en el orden en que aparecen en su lista.
}

\begin{enumerate}
  \item (2 décimas) Dibuje el grafo que representan estas listas.
  \item (1 décimas) ¿El grafo es dirigido? ¿Por qué?
  \item (1 décimas) ¿El grafo es ponderado? ¿Por qué?
  \item (6 décimas) Ejecute DFS recursivo desde \(v_1\). Diligencie la tabla a continuación con el estado de las estructuras en el instante justo después de que se invoca un llamado al procedimiento recursivo de DFS. Liste los llamados a la función en la pila de llamados en el orden en el que ingresaron a la pila, de izquierda a derecha.

\begin{center}
\begin{tabularx}{\textwidth}{|>{\centering\arraybackslash}X|>{\centering\arraybackslash}X|>{\centering\arraybackslash}X|}
\hline
\textbf{Visitado} & \textbf{Pila de llamados} & \textbf{Predecesores} \\
\hline
$v_1$ \rule{0pt}{0.9cm} & & \\
\hline
\rule{0pt}{0.9cm} & & \\
\hline
\rule{0pt}{0.9cm} & & \\
\hline
\rule{0pt}{0.9cm} & & \\
\hline
\rule{0pt}{0.9cm} & & \\
\hline
\rule{0pt}{0.9cm} & & \\
\hline
\rule{0pt}{0.9cm} & & \\
\hline
\end{tabularx}
\end{center}
\end{enumerate}

\bigskip
